\appendix
\cyrchapappendix

\chapter{Численный пример}

Численный пример расчета в настоящем разделе приведен для корпуса по типу контейнеровоза S175 (Вариант 1)

\section{Исходные данные}

Рассмотрен вариант корпуса по типу контейнеровоза S175 (Вариант 1) со следующими характеристиками:

\begin{itemize}
    \item Длина по КВЛ $L_{\text{ПП}} = 200{,}0\,\text{м}$;
    \item Отношение $L/B = 6{,}890$;
    \item Отношение $B/T = 2{,}674$.   
\end{itemize}

Вычислим главные размерения и характеристики, необходимые для построения вспомогательного теоретического чертежа и дальнейших расчетов.


Ширина судна по КВЛ составляет

\begin{equation*}
    B = \frac{L}{L/B} = \frac{200{,}0}{6{,}890} = 29{,}03\,\text{м}.
\end{equation*}

Осадка судна по КВЛ составляет

\begin{equation*}
    T = \frac{B}{B/T} = \frac{29{,}03}{2{,}674} = 10{,}85\,\text{м}.
\end{equation*}

Расстояние между шпангоутами составляет

\begin{equation*}
    \Delta L = L / 10 = 200{,}0 / 10 = 20{,}0\,\text{м}.
\end{equation*}

Конструктивная ватерлиния соответствует пятой ватерлинии, поэтому расстояние между ватерлиниями составляет

\begin{equation*}
    \Delta T = T / 5 = 10{,}85 / 5 = 2{,}17\,\text{м}.
\end{equation*}

Размерные ординаты корпуса приведены в таблице \ref{add-input}.

Расчет коэффициентов учета истинного притыкания ватерлиний в носу и корме приведен в таблице.
Для расчета использованы значения безразмерных величин $\bar{x}_{\text{Ф}}$ и $\bar{x}_{\text{А}}$.

\textit{Примечания — при наличии однотипных расчетов, например, для нескольких ватерлиний, в численном примере приводится расчет только одного случая.}

% Размерные ординаты корпуса по типу контейнеровоза S175
\begin{table}[H]
    \centering
    \caption{Размерные ординаты корпуса по типу контейнеровоза S175}
    \scriptsize
    \begin{tikzpicture}
        \matrix (mat) [addenum,
        column 1/.append style={nodes={text width=0.8cm}},
        column 2/.append style={nodes={text width=0.8cm}},
        column 3/.append style={nodes={text width=0.8cm}},
        column 4/.append style={nodes={text width=0.8cm}},
        column 5/.append style={nodes={text width=0.8cm}},
        column 6/.append style={nodes={text width=0.8cm}},
        column 7/.append style={nodes={text width=0.8cm}},
        column 8/.append style={nodes={text width=0.8cm}},
        column 9/.append style={nodes={text width=0.8cm}},
        column 10/.append style={nodes={text width=0.8cm}},
        column 11/.append style={nodes={text width=0.8cm}},
        column 12/.append style={nodes={text width=0.8cm}}
        ] {
        ВЛ & 0     & 1     & 2     & 3     & 4     & 5     & 6     & 7     & 8     & 9     & 10  \\
        0  &  --- & 0,106 & 1,358 &4,394 & 7,901 & 10,407 & 8,422 & 4,177 & 0,892 & 0,307 & --- \\
        1  &  1,502 & 2,608 & 4,696 &8,439 & 12,225 & 14,250 & 13,204 & 9,709 & 5,441 & 2,068 & --- \\
        2  &  1,360 & 2,766 & 5,791 &9,862 & 13,364 & 14,514 & 14,176 & 11,722 & 7,360 & 2,756 & --- \\
        3  &  0,748 & 2,793 & 6,531 &10,753 & 13,830 & 14,514 & 14,461 & 12,908 & 9,065 & 3,745 & --- \\
        4  &  0,303 & 3,037 & 7,220 &11,456 & 14,067 & 14,514 & 14,514 & 13,662 & 10,703 & 5,551 & --- \\
        5  &  0,006 & 3,647 & 8,057 &12,064 & 14,228 & 14,514 & 14,514 & 14,147 & 12,191 & 8,094 & 1,795 \\
        6  &  0,392 & 4,707 & 9,085 &12,615 & 14,365 & 14,514 & 14,514 & 14,426 & 13,303 & 10,186 & 4,419 \\
        7  &  1,154 & 6,193 & 10,294 &13,194 & 14,467 & 14,514 & 14,514 & 14,514 & 14,130 & 11,720 & 6,161 \\
        П  &  6,652 & 12,292 & 13,854 &14,270 & 14,514 & 14,514 & 14,514 & 14,514 & 14,514 & 13,487 & 8,556 \\
        $z$  &  22,443 & 21,634 & 20,173 &18,971 & 18,971 & 18,971 & 18,971 & 18,971 & 18,971 & 18,971 & 18,971 \\
        $z_1$  &  20,054 & 13,220 & 2,454 &0,051 & 0,000 & 0,000 & 0,000 & 0,080 & 1,619 & 7,963 & 13,551 \\
        $z_2$  &  --- & 18,874 & 14,090 &3,989 & 0,419 & 0,000 & 0,000 & 0,100 & 2,158 & 7,324 & 12,499 \\
        };
    \end{tikzpicture}
    \label{add-input}
\end{table}

% Пример расчета коэффициентов учета истинного притыкания ватерлиний
\begin{table}[H]
    \centering
    \scriptsize
    \caption{Пример расчета коэффициентов учета истинного притыкания ватерлиний}
    \caption*{\textit{
        Схема расчета коэффициентов приведена на рисунке \ref{fig:K-values}.
    }}
    \begin{tikzpicture}
        \matrix (mat) [addenum,
        column 1/.append style={nodes={text width=1.0cm}},
        column 2/.append style={nodes={text width=1.5cm}},
        column 3/.append style={nodes={text width=1.5cm}},
        column 4/.append style={nodes={text width=1.0cm}},
        column 5/.append style={nodes={text width=1.0cm}},
        column 6/.append style={nodes={text width=1.5cm}},
        column 7/.append style={nodes={text width=1.5cm}}
        ] {
        $i_{\text{ВЛ}}$ & $\bar{x}_{\text{Ф}}$ & $\bar{x}_{\text{А}}$ & $i_{\text{Н}}$ & $i_{\text{К}}$ & $K_{\text{Ф}}$ & $K_{\text{А}}$\\
        \scriptsize{I}   &    \scriptsize{II}   &   \scriptsize{III}   &   \scriptsize{IV} &   \scriptsize{V} &   \scriptsize{VI} &   \scriptsize{VII} \\
        0  &  -0,0686  &  0,5609  &  1  &  9  &  0,9657  &  0,7195   \\
        1  &  0,0964  &  0,5185  &  0  &  9  &  0,5482  &  0,7408   \\
        2  &  0,0844  &  0,4883  &  0  &  9  &  0,5422  &  0,7558   \\
        3  &  0,0379  &  0,4852  &  0  &  9  &  0,5189  &  0,7574   \\
        4  &  0,0116  &  0,4591  &  0  &  9  &  0,5058  &  0,7705   \\
        5  &  0,0000  &  0,0000  &  0  &  10  &  0,5000  &  0,5000   \\
        6  &  0,0100  &  -0,1297  &  0  &  10  &  0,5050  &  0,5649   \\
        7  &  0,0393  &  -0,1687  &  0  &  10  &  0,5196  &  0,5844   \\
        };
    \end{tikzpicture}
\end{table}

\newpage

\section{Построение кривых элементов теоретического чертежа}

По размерным значениям ординат для каждой ватерлинии выполняется расчет площади ватерлинии, статического момента площади ватерлинии и моментов инерции площади ватерлинии относительно продольной и поперечной осей.
Пример расчета для четвертой ватерлинии приведен в таблицах ниже.

% Пример расчета элементов площади ватерлинии (площади и статического момента)
\begin{table}[H]
    \centering
    \scriptsize
    \caption{Пример расчета площади ватерлинии и статического момента площади ватерлинии от осадки}
    \caption*{\textit{
        Столбец III условно назван $S$ --- значение площади ватерлинии рассчитывается в последней строке столбца.
        Столбец IV условно назван $M_y$ --- значение статического момента площади ватерлинии рассчитывается в последней строке столбца.
        Пример расчета приведен для четвертой ватерлинии.
        Коэффициент $K_{\text{Ф}} = 0,506$ , а коэффициент $K_{\text{А}} = 0,770$.
        Схема расчета приведена на рисунке \ref{tab:waterline-elements}.
    }}
    \begin{tikzpicture}
        \matrix (mat) [addenum,
        column 1/.append style={nodes={text width=1.0cm}},
        column 2/.append style={nodes={text width=2.5cm}},
        column 3/.append style={nodes={text width=2.5cm}},
        column 4/.append style={nodes={text width=2.5cm}}
        ] {
        $i_{\text{Ш}}$ & $y$ & $S$ & $M_y$  \\
        \scriptsize{I}   &    \scriptsize{II}   &   \scriptsize{III}   &   \scriptsize{IV} \\
        0   &   0,303   &   0,153   &   0,77   \\
        1   &   3,037   &   3,037   &   12,15   \\
        2   &   7,220   &   7,220   &   21,66   \\
        3   &   11,456   &   11,456   &   22,91   \\
        4   &   14,067   &   14,067   &   14,07   \\
        5   &   14,514   &   14,514   &   0,00   \\
        6   &   14,514   &   14,514   &   -14,51   \\
        7   &   13,662   &   13,662   &   -27,32   \\
        8   &   10,703   &   10,703   &   -32,11   \\
        9   &   5,551   &   4,277   &   -17,11   \\
        10   &      &      &      \\
        ---   &   ---   &   3744,2   &   -15599,6   \\
        };
    \end{tikzpicture}
\end{table}

% Пример расчета элементов площади ватерлинии (моментов инерции относительно продольной и поперечной осей)
\begin{table}[H]
    \centering
    \scriptsize
    \caption{Пример расчета моментов инерции площади ватерлинии относительно продольной и поперечной осей}
    \caption*{\textit{
        Столбец III условно назван $I_x$ --- значение момента инерции площади ватерлинии рассчитывается в последней строке столбца.
        Столбец IV условно назван $I_y$ --- значение  момента инерции площади ватерлинии рассчитывается в последней строке столбца.
        Пример расчета приведен для четвертой ватерлинии.
        Коэффициент $K_{\text{Ф}} = 0,506$ , а коэффициент $K_{\text{А}} = 0,770$.
        Схема расчета приведена на рисунке \ref{tab:inertia-elements}.
    }}
    \begin{tikzpicture}
        \matrix (mat) [addenum,
        column 1/.append style={nodes={text width=1.0cm}},
        column 2/.append style={nodes={text width=2.5cm}},
        column 3/.append style={nodes={text width=2.5cm}},
        column 4/.append style={nodes={text width=2.5cm}}
        ] {
        $i_{\text{Ш}}$ & $y$ & $I_x$ & $I_y$  \\
        \scriptsize{I}   &    \scriptsize{II}   &   \scriptsize{III}   &   \scriptsize{IV} \\
        0   &   0,303   &   0,014   &   3,83   \\
        1   &   3,037   &   28,006   &   48,59   \\
        2   &   7,220   &   376,400   &   64,98   \\
        3   &   11,456   &   1503,677   &   45,83   \\
        4   &   14,067   &   2783,687   &   14,07   \\
        5   &   14,514   &   3057,645   &   0,00   \\
        6   &   14,514   &   3057,645   &   14,51   \\
        7   &   13,662   &   2550,155   &   54,65   \\
        8   &   10,703   &   1225,987   &   96,32   \\
        9   &   5,551   &   131,768   &   68,43   \\
        10   &      &      &      \\
        —   &   —   &   196199,8   &   6579359,4   \\
        };
    \end{tikzpicture}
\end{table}

Аналогично выполняется расчет для остальных ватерлиний.
Полученные величины сводятся в одну таблицу в зависимости от осадки.

% Элементы площади ватерлинии в зависимости от осадки
\begin{table}[H]
    \centering
    \scriptsize
    \caption{Элементы площади ватерлинии в зависимости от осадки}
    \caption*{\textit{
        Схема расчета приведена на рисунке \ref{tab:waterline-elements-results}.
    }}
    \begin{tikzpicture}
        \matrix (mat) [addenum,
        column 1/.append style={nodes={text width=1.0cm}},
        column 2/.append style={nodes={text width=1.4cm}},
        column 3/.append style={nodes={text width=1.4cm}},
        column 4/.append style={nodes={text width=1.4cm}},
        column 5/.append style={nodes={text width=1.4cm}},
        column 6/.append style={nodes={text width=1.4cm}},
        column 7/.append style={nodes={text width=1.4cm}},
        column 8/.append style={nodes={text width=1.4cm}}
        ] {
        $i_{\text{ВЛ}}$ & $z$ & $S$ & $x_f$ & $M_y$ & $I_x$  & $I_y$ & $I_{yf}$ \\
        \scriptsize{I}   &    \scriptsize{II}   &   \scriptsize{III}   &   \scriptsize{IV} &   \scriptsize{V} &   \scriptsize{VI} &   \scriptsize{VII} &   \scriptsize{VIII}\\
        0   &   0,000   &   1515,0   &   0,442   &   670,1   &   31717,2   &   1216477   &   1216180   \\
        1   &   2,171   &   2917,1   &   0,731   &   2133,8   &   117735,5   &   4417208   &   4415648   \\
        2   &   4,343   &   3295,0   &   -0,685   &   -2256,7   &   153254,3   &   5251966   &   5250420   \\
        3   &   6,514   &   3523,2   &   -2,448   &   -8623,9   &   176087,0   &   5809196   &   5788087   \\
        4   &   8,686   &   3744,2   &   -4,166   &   -15599,6   &   196199,8   &   6579359   &   6514366   \\
        5   &   10,857   &   4094,2   &   -7,644   &   -31294,7   &   219988,9   &   8418697   &   8179493   \\
        6   &   13,029   &   4416,4   &   -9,027   &   -39865,6   &   245376,0   &   10307031   &   9947177   \\
        7   &   15,200   &   4709,7   &   -8,715   &   -41045,2   &   271932,0   &   12019817   &   11662104   \\
        };
    \end{tikzpicture}
\end{table}


Для вычисления элементов погруженного объема --- водоизмещения, статических моментов объема, центра тяжести водоизмещения --- расчет выполняется в соответствии со схемой, приведенной на рисунке \ref{tab:volume-elements}

% Расчет объемного водоизмещения и статических моментов объема
\begin{table}[H]
    \centering
    \scriptsize
    \caption{Пример расчета объемного водоизмещения и статических моментов объема}
    \caption*{\textit{
        Схема расчета приведена на рисунке \ref{tab:volume-elements}.
    }}
    \begin{tikzpicture}
        \matrix (mat) [addenum,
        column 1/.append style={nodes={text width=1.0cm}},
        column 2/.append style={nodes={text width=1.4cm}},
        column 3/.append style={nodes={text width=1.4cm}},
        column 4/.append style={nodes={text width=1.4cm}},
        column 5/.append style={nodes={text width=1.4cm}},
        column 6/.append style={nodes={text width=1.4cm}},
        column 7/.append style={nodes={text width=1.4cm}},
        column 8/.append style={nodes={text width=1.4cm}}
        ] {
        $i_{\text{ВЛ}}$ & $S$ & $M_y$ & $V$ & $M_{yz}$ & $M_{xy}$  & $x_c$ & $z_c$ \\
        \scriptsize{I}   &    \scriptsize{II}   &   \scriptsize{III}   &   \scriptsize{IV} &   \scriptsize{V} &   \scriptsize{VI} &   \scriptsize{VII} &   \scriptsize{VIII}\\
        0   &   1515,0   &   670,1   &   0   &   0   &   0   &   ---   &   ---   \\
        1   &   2917,1   &   2133,8   &   4812   &   3044   &   6877   &   0,633   &   1,429   \\
        2   &   3295,0   &   -2256,7   &   11557   &   2911   &   29291   &   0,252   &   2,535   \\
        3   &   3523,2   &   -8623,9   &   18959   &   -8902   &   69745   &   -0,470   &   3,679   \\
        4   &   3744,2   &   -15599,6   &   26849   &   -35202   &   129972   &   -1,311   &   4,841   \\
        5   &   4094,2   &   -31294,7   &   35360   &   -86116   &   213542   &   -2,435   &   6,039   \\
        6   &   4416,4   &   -39865,6   &   44600   &   -163376   &   324276   &   -3,663   &   7,271   \\
        7   &   4709,7   &   -41045,2   &   54508   &   -251222   &   464470   &   -4,609   &   8,521   \\
        };
    \end{tikzpicture}
\end{table}

Расчет метацентрических радиусов и возвышения метацентра выполняется в соответствии со схемой, приведенной на рисунке \ref{tab:metacentric-radii}.

% Расчет метацентрических радиусов и возвышения метацентра
\begin{table}[H]
    \centering
    \scriptsize
    \caption{Пример расчета метацентрических радиусов и возвышения метацентра}
    \caption*{\textit{
        Схема расчета приведена на рисунке \ref{tab:metacentric-radii}.
    }}
    \begin{tikzpicture}
        \matrix (mat) [addenum,
        column 1/.append style={nodes={text width=1.0cm}},
        column 2/.append style={nodes={text width=1.4cm}},
        column 3/.append style={nodes={text width=1.4cm}},
        column 4/.append style={nodes={text width=1.4cm}},
        column 5/.append style={nodes={text width=1.4cm}},
        column 6/.append style={nodes={text width=1.4cm}},
        column 7/.append style={nodes={text width=1.4cm}},
        column 8/.append style={nodes={text width=1.4cm}}
        ] {
        $i_{\text{ВЛ}}$ & $V$ & $I_x$ & $I_{yf}$ & $z_c$ & $r$  & $R$ & $z_m$ \\
        \scriptsize{I}   &    \scriptsize{II}   &   \scriptsize{III}   &   \scriptsize{IV} &   \scriptsize{V} &   \scriptsize{VI} &   \scriptsize{VII} &   \scriptsize{VIII}\\
        0   &   0   &   31717   &   1216180   &   ---   &   ---   &   ---   &   ---   \\
        1   &   4812   &   117736   &   4415648   &   1,429   &   24,47   &   917,6   &   25,90   \\
        2   &   11557   &   153254   &   5250420   &   2,535   &   13,26   &   454,3   &   15,80   \\
        3   &   18959   &   176087   &   5788087   &   3,679   &   9,29   &   305,3   &   12,97   \\
        4   &   26849   &   196200   &   6514366   &   4,841   &   7,31   &   242,6   &   12,15   \\
        5   &   35360   &   219989   &   8179493   &   6,039   &   6,22   &   231,3   &   12,26   \\
        6   &   44600   &   245376   &   9947177   &   7,271   &   5,50   &   223,0   &   12,77   \\
        7   &   54508   &   271932   &   11662104   &   8,521   &   4,99   &   214,0   &   13,51   \\
        };
    \end{tikzpicture}
\end{table}

Результат расчета кривых элементов теоретического чертежа приведен на рисунке.

\textit{
    Примечание --- В численном примере графики кривых элементов теоретического чертежа приведены в упрощенном виде (в виде диаграммы, построенной с применением табличного редактора) в связи с форматом пособия.
    В курсовой работе диаграммы должны быть построены  в CAD-системе на основе расчетных данных, приведенных в таблицах, с соблюдением всех требований, указанных в разделе 2.
}

\begin{figure}[H]
    \centerline{\incfig{add_01}}
    \caption{Кривые элементов теоретического чертежа}
    \caption*{\textit{}}
\end{figure}

\newpage
\section{Расчет масштаба Бонжана}

Для всех одиннадцати шпангоутов выполняется расчет зависимости площади шпангоута от осадки по схеме, приведенной на рисунке \ref{tab:station-area-bonjean}.
Ниже приведен пример расчета для пятого шпангоута.

% Расчет зависимости площади шпангоута от осадки
\begin{table}[H]
    \centering
    \scriptsize
    \caption{Пример расчета зависимости площади шпангоута от осадки}
    \caption*{\textit{
        Схема расчета приведена на рисунке \ref{tab:station-area-bonjean}.
        Расчет выполнен для пятого шпангоута.
    }}
    \begin{tikzpicture}
        \matrix (mat) [addenum,
        column 1/.append style={nodes={text width=1.0cm}},
        column 2/.append style={nodes={text width=2.0cm}},
        column 3/.append style={nodes={text width=2.0cm}},
        column 4/.append style={nodes={text width=2.0cm}}
        ] {
        $i_{\text{ВЛ}}$ & $z$ & $y$ & $\Omega $ \\
        \scriptsize{I}   &    \scriptsize{II}   &   \scriptsize{III}   &   \scriptsize{IV} \\
        0  &  0,000  &  10,407  &  0,000   \\
        1  &  2,171  &  14,250  &  49,386   \\
        2  &  4,343  &  14,514  &  111,845   \\
        3  &  6,514  &  14,514  &  174,879   \\
        4  &  8,686  &  14,514  &  237,912   \\
        5  &  10,857  &  14,514  &  300,946   \\
        6  &  13,029  &  14,514  &  363,979   \\
        7  &  15,200  &  14,514  &  427,013   \\
        ---  &  18,971  &  14,514  &  536,492   \\
        ---  &  19,552  &  ---  &  547,727   \\
        };
    \end{tikzpicture}
\end{table}

Полученный по результатам расчета масштаб Бонжана приведен на рисунке.

% Масштаб Бонжана
\begin{figure}[H]
    \centerline{\incfig{add_02}}
    \caption{Масштаб Бонжана}
    \caption*{\textit{}}
\end{figure}

В численном примере принята средняя осадка $T = 13{,}029\,\text{м}$, а угол дифферента $\psi = 1,0\,\text{\textdegree}$.

Тогда осадка на носовом перпендикуляре составляет

\begin{equation*}
    T_{\text{Н}} = T + \psi (L/2 - x_f) = 13{,}029 + 0{,}01745 (200 / 2 - (-9{,}027)) = 14{,}93\,\text{м}.
\end{equation*}

Осадка на кормовом перпендикуляре составляет

\begin{equation*}
    T_{\text{К}} = T - \psi (L/2 + x_f) = 13{,}029 - 0{,}01745 (200 / 2 + (-9{,}027)) = 11{,}44\,\text{м}.
\end{equation*}

При этих значениях осадки на перпендикулярах коэффициенты учета истинного притыкания ватерлиний составляют $K_{\text{Н}} = 0{,}518$ и $K_{\text{К}} = 0{,}486$.

По значениям осадки на перпендикулярах на масштабе Бонжана проводится наклонная прямая --- ватерлиния судна, плавающего с дифферентом.
Для расчета водоизмещения и абсциссы центра величины с масштаба Бонжана снимаются значения площади погруженной части каждого шпангоута.
В соответствии со схемой расчета, представленной на рисунке \ref{fig:volume-elements-bonjean}, выполняется расчет водоизмещения и абсциссы центра величины судна.

% Пример расчета водоизмещения и абсциссы центра величины судна по масштабу Бонжана
\begin{table}[H]
    \centering
    \scriptsize
    \caption{Пример расчета водоизмещения и абсциссы центра величины судна по масштабу Бонжана}
    \caption*{\textit{
        Столбец III условно назван $V$ --- значение водоизмещения рассчитывается в последней строке столбца.
        Столбец IV условно назван $x_c$ --- значение абсциссы центра величины рассчитывается в последней строке столбца.
        Схема расчета приведена на рисунке \ref{fig:volume-elements-bonjean}.
    }}
    \begin{tikzpicture}
        \matrix (mat) [addenum,
        column 1/.append style={nodes={text width=1.0cm}},
        column 2/.append style={nodes={text width=2.0cm}},
        column 3/.append style={nodes={text width=2.0cm}},
        column 4/.append style={nodes={text width=2.0cm}}
        ] {
        $i_{\text{ШП}}$ & $\Omega$ & $V$ & $x_c $ \\
        \scriptsize{I}   &    \scriptsize{II}   &   \scriptsize{III}   &   \scriptsize{IV} \\
        0  &  20,18  &  10,45  &  52,23   \\
        1  &  90,78  &  90,78  &  363,14   \\
        2  &  185,08  &  185,08  &  555,24   \\
        3  &  286,58  &  286,58  &  573,15   \\
        4  &  356,71  &  356,71  &  356,71   \\
        5  &  372,63  &  372,63  &  0,00   \\
        6  &  352,10  &  352,10  &  -352,10   \\
        7  &  294,89  &  294,89  &  -589,77   \\
        8  &  202,08  &  202,08  &  -606,25   \\
        9  &  95,60  &  95,60  &  -382,41   \\
        10  &  2,70  &  1,31  &  -6,57   \\
        ---  &  ---  &  44960,00  &  -0,33   \\
        };
    \end{tikzpicture}
\end{table}

В результате расчета получено значение водоизмещения и абсциссы центра величины
\begin{equation*}
    \begin{split}
    V = 44960\,\text{м}^3, \\
    x_c = -0{,}33\,\text{м}.
    \end{split}
\end{equation*}

\newpage

\section{Расчет остойчивости на больших углах крена}

Расчет метацентрических радиусов при накренении судна выполнен для осадки по четвертую ватерлинию.
Исходными данными для графо-аналитического расчета являются значения объемного водоизмещения, метацентрического радиуса аппликаты центра величины и коэффициентов учета притыкания ватерлинии ненакрененного судна для расчетной вателинии

\begin{equation*}
    \begin{split}
        V = 26849\,\text{м}^3,\\
        r = 7{,}307\,\text{м}, \\
        z_c = 4{,}841\,\text{м}, \\
        K_{\text{Ф}} = 0{,}5058, \\
        K_{\text{А}} = 0{,}7705.
    \end{split}
\end{equation*}

В соответствии со схемой расчета на вспомогательном теоретическом чертежа проводится ватерлиния, наклоненная на 15\textdegree.
Для каждого шпангоута с чертежа снимаются значения ординаты борта, входящего в воду (в нашем случае --- правого) и ординат борта, выходящего из воды, относительно точки наклона.
С использованием снятых ординат выполняется расчет метацентрического радиуса и поправки для обеспечения равнообъемности.

% Расчет метацентрического радиуса и поправки для обеспечения равнообъемности при накренении
\begin{table}[H]
    \centering
    \scriptsize
    \caption{Пример расчета метацентрического радиуса и поправки для обеспечения равнообъемности при накренении судна на 15\textdegree}
    \caption*{\textit{
        Столбец IV условно назван $S_{\Theta}$ --- значение водоизмещения рассчитывается в последней строке столбца.
        Столбец V условно назван $M_{y \Theta}$ --- значение статического момента площади рассчитывается в последней строке столбца.
        Столбец VI условно назван $I'_{x \Theta}$ --- значение момента инерции площади рассчитывается в последней строке столбца.
        Схема расчета приведена на рисунке \ref{tab:kryloff-method-II}.
    }}
    \begin{tikzpicture}
        \matrix (mat) [addenum,
        column 1/.append style={nodes={text width=0.8cm}},
        column 2/.append style={nodes={text width=1.6cm}},
        column 3/.append style={nodes={text width=1.6cm}},
        column 4/.append style={nodes={text width=1.6cm}},
        column 5/.append style={nodes={text width=2.0cm}},
        column 6/.append style={nodes={text width=2.0cm}}
        ] {
        $i_{\text{ШП}}$ & $y_{\text{пр}}$ & $y_{\text{лб}}$ & $S_{\Theta}$ & $M_{y \Theta}$ &  $I'_{x \Theta}$ \\
        \scriptsize{I}   &    \scriptsize{II}   &   \scriptsize{III}   &   \scriptsize{IV} & \scriptsize{V} & \scriptsize{VI}\\
        0  &  0,297  &  -0,332  &  0,32  &  -0,01  &  0,03   \\
        1  &  3,343  &  -3,016  &  6,36  &  2,08  &  64,78   \\
        2  &  8,334  &  -6,886  &  15,22  &  22,04  &  905,41   \\
        3  &  12,789  &  -10,889  &  23,68  &  45,00  &  3383,17   \\
        4  &  14,842  &  -14,032  &  28,87  &  23,37  &  6031,98   \\
        5  &  15,026  &  -14,789  &  29,82  &  7,07  &  6627,54   \\
        6  &  15,026  &  -14,789  &  29,82  &  7,07  &  6627,54   \\
        7  &  14,895  &  -12,794  &  27,69  &  58,18  &  5399,15   \\
        8  &  13,369  &  -9,215  &  22,58  &  93,81  &  3171,84   \\
        9  &  8,375  &  -4,568  &  9,97  &  37,96  &  526,05   \\
        10  &  ---  &  ---  &  ---  &  ---  &  ---   \\
        ---  &  ---  &  ---  &  3886,5  &  2966  &  218250   \\
        };
    \end{tikzpicture}
\end{table}

Ордината центра тяжести площади наклоненной ватерлинии 
\begin{equation*}
    y_{f 15} = \frac{M_{y 15}}{S_{15}} = \frac{2966}{3886{,}5} = 0{,}763\,\text{м}.
\end{equation*}

Момент инерции площади наклоненной ватерлинии относительно оси, проходящей через центр тяжести площади
\begin{equation*}
    I_{x 15} = I'_{x 15} - S_{15} y_{f 15}^2 = 218250 - 3886{,}5 \cdot 0{,}763^2 = 215990\,\text{м}^4.
\end{equation*}

Тогда метацентрический радиус при накренении судна на 15\textdegree составляет
\begin{equation*}
    r_{15} = \frac{I_{x 15}}{V} = \frac{215990}{26849} = 8{,}04\,\text{м}.
\end{equation*}

С учетом ординаты центра тяжести площади ватерлинии на вспомогательном теоретическом чертеже проводится истинная равнообъемная ватерлиния, от которой выполняется аналогичный расчет при накренении на 30\textdegree.

После выполнения расчетов для всех углов накренения (15\textdegree, 30\textdegree, 45\textdegree, 60\textdegree) выполняется построение кривой метацентрических радиусов в зависимости от угла накренения для нескольких значений водоизмещения судна (в примере расчет выполнен для ватерлиний с третьей по седьмую).

% Зависимость метацентрического радиуса от угла крена
\begin{figure}[H]
    \centerline{\incfig{add_03}}
    \caption{Зависимость метацентрического радиуса от угла крена}
\end{figure}

По полученным значениям метацентрических радиусов выполняется расчет плеча остойчивости формы.
Пример расчета для одного из значений водоизмещения судна приведен в таблице.

% Расчет плеча остойчивости формы
\begin{table}[H]
    \centering
    \footnotesize
    \caption{Пример расчета плеча остойчивости формы}
    \caption*{\textit{
        Схема расчета приведена на рисунке \ref{tab:kryloff-method-II}.
    }}
    \begin{tikzpicture}
        \matrix (mat) [addenum,
        column 1/.append style={nodes={text width=1.0cm}},
        column 2/.append style={nodes={text width=2.0cm}},
        column 3/.append style={nodes={text width=2.0cm}},
        column 4/.append style={nodes={text width=2.0cm}},
        column 5/.append style={nodes={text width=2.0cm}}
        ] {
        $\Theta$ & $r$ & $y_{\Theta}$ & $z_{\Theta}- z_c$ & $l_{\text{Ф}}$  \\
        \scriptsize{I}   &    \scriptsize{II}   &   \scriptsize{III}   &   \scriptsize{IV} & \scriptsize{V} \\
        0  &  7,307  &  0,000  &  0,000  &  0,000   \\
        15  &  8,044  &  1,974  &  0,273  &  1,977   \\
        30  &  9,726  &  4,093  &  1,182  &  4,136   \\
        45  &  8,548  &  5,987  &  2,609  &  6,079   \\
        60  &  6,278  &  7,189  &  4,112  &  7,156   \\
        };
    \end{tikzpicture}
\end{table}

По результатам расчета плеча остойчивости формы для всех значений водоизмещения строится диаграмма зависимости плеча остойчивости формы от угла крена для нескольких значений водоизмещения судна.

% Зависимость плеча остойчивости формы от угла крена
\begin{figure}[H]
    \centerline{\incfig{add_04}}
    \caption{Зависимость плеча остойчивости формы от угла крена}
\end{figure}

Интерполяционные кривые плеч остойчивости формы получаются перестроением кривых плеч остойчивости формы в оси <<Водоизмещение>> --- <<Плечо остойчивости формы>>.
Несколько кривых соответствуют нескольким значениям угла крена (15\textdegree, 30\textdegree, 45\textdegree, 60\textdegree).

% Интерполяционные кривые плеч остойчивости формы
\begin{figure}[H]
    \centerline{\incfig{add_05}}
    \caption{Интерполяционные кривые плеч остойчивости формы}
\end{figure}

Построение диаграммы статической и динамической остойчивости выполнено для водоизмещения судна $V = 40000\,\text{м}^3$ и аппликаты центра тяжести $z_g = 12{,}0\,\text{м}$.
Аппликата центра величины при посадке на ровный киль составляет $z_c = 6{,}67\, \text{м}$.
Тогда
\begin{equation*}
    a = z_g - z_c = 12{,}0 - 6{,}67 = 5{,}33\,\text{м}.
\end{equation*}

Расчет плеч статической и динамической остойчивости выполняется в соответствии со схемой, приведенной на рисунке \ref{fig:stability-curve-calculation}.

% Расчет плеча статической и динамической остойчивости
\begin{table}[H]
    \centering
    \footnotesize
    \caption{Пример расчета диаграммы статической и динамической остойчивости}
    \caption*{\textit{
        В столбце IV выполняется предварительное суммирование для расчета плеча динамической остойчивости.
        В столбце V значения из столбца IV домножаются на $0,5 \Delta\Theta $ (в радианах).
        Схема расчета приведена на рисунке \ref{fig:stability-curve-calculation}.
    }}
    \begin{tikzpicture}
        \matrix (mat) [addenum,
        column 1/.append style={nodes={text width=1.0cm}},
        column 2/.append style={nodes={text width=2.0cm}},
        column 3/.append style={nodes={text width=2.0cm}},
        column 4/.append style={nodes={text width=2.0cm}},
        column 5/.append style={nodes={text width=2.0cm}}
        ] {
        $\Theta$ & $l_{\text{Ф}}$ & $l_{\Theta}$ & --- & $d$  \\
        \scriptsize{I}   &    \scriptsize{II}   &   \scriptsize{III}   &   \scriptsize{IV} & \scriptsize{V} \\
        0  &  0,000  &  0,000  &  0,000  &  0,000   \\
        15  &  1,557  &  0,177  &  0,457  &  0,020   \\
        30  &  3,309  &  0,642  &  2,786  &  0,122   \\
        45  &  4,729  &  0,958  &  7,849  &  0,342   \\
        60  &  5,289  &  0,671  &  12,967  &  0,566   \\
        };
    \end{tikzpicture}
\end{table}

Полученная диаграмма статической и динамической остойчивости приведена на рисунке.

\begin{figure}[H]
    \centerline{\incfig{add_06}}
    \caption{Диаграмма статической и динамической остойчивости}
    \caption*{\textit{
        Для водоизмещения $V = 40000\,\text{м}^3$ и аппликаты центра тяжести $z_g = 12{,}0\,\text{м}$.
    }}
\end{figure}

\section{Расчет непотопляемости судна}

Расчет непотопляемости судна выполнен для следующего случая нагрузки: объемное водоизмещение $V = 35360\,\text{м}^3$, абсцисса центра тяжести $x_g = -2{,}435\,\text{м}$, аппликата центра тяжести $z_g = 12{,}0\,\text{м}$, осадка по ватерлинии $T = 10{,}857\,\text{м}$.
В исходном положении судно сидит на ровный киль.

Абсцисса и аппликата центра величины составляют $x_c = -2{,}435\,\text{м}$ и $z_c = 6{,}039\,\text{м}$ соответственно.
Поперечный и продольный метацентрические радиусы составляют $r = 6{,}22\,\text{м}$ и $R = 231{,}3\,\text{м}$ соответственно.

Площадь ватерлинии составляет $S = 4094{,}2\,\text{м}^2$.
Абсцисса центра тяжести ватерлинии составляет $x_f = -7{,}644\,\text{м}$.

Центральные моменты инерции $I_x = 219990\,\text{м}^4$ и $I_y = 8179493\,\text{м}^4$.

Предельная линия погружения принята проходящей на 0,3 м ниже линии верхней палубы у борта --- $18{,}67\,\text{м}$.

Для проведения касательный к предельной линии погружения вычисляется величина, откладываемая на носовом и кормовом перпендикулярах
\begin{equation*}
    z_0 = 1{,}2 (2 \cdot T_{\text{КВЛ}} - H) = 1{,}2 (2 \cdot 10{,}857 - 18{,}97) = 3{,}29\,\text{м}.
\end{equation*}

В численном примере на масштаб Бонжана нанесено семь предельных ватерлиний.
Значения площади погруженной части каждого шпангоута для каждой из предельных ватерлиний сняты с масштаба Бонжана и используются для вычисления объема и положения центра тяжести отсеков, при затоплении которых судно погрузится до предельной линии.
В таблице ниже приведен пример расчета для одной из предельных ватерлиний.

% Расчет водоизмещения и статического момента
\begin{table}[H]
    \centering
    \footnotesize
    \caption{Пример расчета водоизмещения }
    \caption*{\textit{
        В столбце II приведены снятые с масштаба Бонжана значения.
        Столбец IV условно назван $V$ --- значение водоизмещения рассчитывается в последней строке столбца.
        Столбец IV условно назван $M$ --- значение статического момента  рассчитывается в последней строке столбца.
        Схема расчета приведена на рисунке \ref{fig:full-volume}.
    }}
    \begin{tikzpicture}
        \matrix (mat) [addenum,
        column 1/.append style={nodes={text width=1.0cm}},
        column 2/.append style={nodes={text width=2.0cm}},
        column 3/.append style={nodes={text width=2.0cm}},
        column 4/.append style={nodes={text width=2.0cm}}
        ] {
        $i_{\text{ШП}}$ &  $\Omega$ & $V$ & $M$ \\
        \scriptsize{I}   &    \scriptsize{II}   &   \scriptsize{III}   &   \scriptsize{IV} \\
        0  &  76,99  &  40,00  &  200,02   \\
        1  &  217,75  &  217,75  &  871,00   \\
        2  &  338,58  &  338,58  &  1015,73   \\
        3  &  447,24  &  447,24  &  894,49   \\
        4  &  523,30  &  523,30  &  523,30   \\
        5  &  534,53  &  534,53  &  0,00   \\
        6  &  505,31  &  505,31  &  -505,31   \\
        7  &  436,44  &  436,44  &  -872,87   \\
        8  &  324,17  &  324,17  &  -972,52   \\
        9  &  180,66  &  180,66  &  -722,63   \\
        10  &  34,29  &  20,11  &  -100,55   \\
        ---  &  ---  &  71362  &  132260   \\
        };
    \end{tikzpicture}
\end{table}

Тогда объем отсека при затоплении которого судно погрузится до предельной линии составляет
\begin{equation*}
    v_2 = V_{2-2} - V_0 = 71362 - 35360 = 36002\,\text{м}^3.
\end{equation*}

Абсцисса центра тяжести отсека составит
\begin{equation*}
    x_2 = \frac{M_{2-2} - M_0}{v_2} = \frac{132260 - (-86116)}{36002} = 6{,}07\,\text{м}.
\end{equation*}

По полученным для семи предельный ватерлиний  строится кривая предельных отсеков.

% Кривая предельных объемов
\begin{figure}[H]
    \centerline{\incfig{add_07}}
    \caption{Кривая предельных объемов}
\end{figure}

Пример расчета интегральной кривой от строевой по шпангоутам приведен в таблице и на рисунке.

% Интегральная кривая от строевой по шпангоутам
\begin{figure}[H]
    \centerline{\incfig{add_08}}
    \caption{Интегральная кривая от строевой по шпангоутам}
\end{figure}

В соответствии со схемой расчета, представленной на рисунке \ref{fig:max-volume-curve}, выполнен расчет кривой предельных длин отсеков.
Построения приведены на рисунке.

\begin{figure}[H]
    \centerline{\incfig{add_09}}
    \caption{К определению предельных длин отсеков}
    \caption*{\textit{
    }}
\end{figure}

Полученная кривая предельных длин отсеков приведена на рисунке.

\begin{figure}[H]
    \centerline{\incfig{add_10}}
    \caption{Кривая предельных длин отсеков}
    \caption*{\textit{
    }}
\end{figure}

Определим посадку и начальную остойчивость поврежденного судна.

В численном примере принимаем, что затапливаемый отсек располагается между вторым и четвертым шпангоутами.
С учетом полученной кривой предельных длин отсеков, предельная длина отсека составляет $l_{\text{пр}} \approx 40{,}0\,\text{м}$.
Таким образом, учитывая, что расстояние между шпангоутами в численном примере составляет $\Delta L = 20\,\text{м}$, носовая переборка отсека соответствует второму шпангоуту, а кормовая переборка отсека --- четвертому шпангоуту.

В расчете примем, что до затопления судно сидит но ровный киль по пятую ватерлинию. Коэффициент проницаемости для отсека примем $\mu = 1{,}0$.

Площади погруженной части второго, третьего и четвертого шпангоутов составляют (по масштабу Бонжана)

\begin{equation*}
    \Omega_2 = 125{,}70\,\text{м}^2, \quad \Omega_3 = 211{,}67\,\text{м}^2, \quad \Omega_4 = 280{,}33\,\text{м}^2.
\end{equation*}

Объем отсека можно определить как площадь двух трапеций (между носовой переборкой и третьим шпангоутом и между третьим шпангоутом и кормовой переборкой)

\begin{equation*}
    \begin{split}
    v = \frac{\Omega_2 + \Omega_3}{2} \cdot l + \frac{\Omega_3+\Omega_4}{2} \cdot l = \\
    = \frac{125{,}70 + 211{,}67}{2} \cdot 20 + \frac{211{,}67 + 280{,}33}{2} \cdot 20 = 8293{,}8\,\text{м}^3.
    \end{split}
\end{equation*}

Положение центра тяжести затопленного отсека определим интегрированием строевой по шпангоутам

\begin{equation*}
    x_{\text{р}} = \frac{1}{v} \int_{II}^{I} \Omega x dx , \quad
    z_{\text{р}} = T - \frac{1}{v} \int_{II}^{I} \int_0^T \Omega dz dx
\end{equation*}

Получены следующие значения

\begin{equation*}
    x_{\text{р}} = 37{,}46\,\text{м} , \quad
    z_{\text{р}} = 5{,}97 \,\text{м}
\end{equation*}

Элементы потерянной площади ватерлинии определяются по ординатам действующей ватерлинии --- $y_2 = 8{,}057\,\text{м}$, $y_3 = 12{,}064\,\text{м}$, $y_4 = 14{,}228\,\text{м}$.

Потерянная площадь определяется как

\begin{equation*}
    s_0 = l/2 \cdot (y_2 + 2 \cdot y_3 + y_4) = 40 / 2 \cdot (8{,}057 + 2 \cdot 12{,}064 + 14{,}288) = 928{,}24\,\text{м}^2
\end{equation*}

Абсцисса центра тяжести потерянной площади определяется как

\begin{equation*}
    \begin{split}
    x_s = \frac{1}{s_0} \int_{II}^{I} y x dx = \\
    = \frac{1}{s_0} (l/2 \cdot(x_2 y_2 + 2 \cdot x_3  y_3 + x_4 y_4)) = \\
    = \frac{1}{928{,}24} (40/2 \cdot(60{,}0 \cdot 8{,}057 + 2 \cdot 40{,}0 \cdot 12{,}064 + 20{,}0 \cdot 14{,}288)) = \\
    = 37{,}34\,\text{м}
    \end{split}
\end{equation*}

Момент инерции потерянной площади относительно продольной оси

\begin{equation*}
    \begin{split}
        i_x = \frac{2}{3} \int_{II}^{I} y^3 dx = \\
        = \frac{2}{3} l/2 (0{,}5 \cdot y_2^3 + y_3^3 + 0{,}5 \cdot y_4^3) = \\
        = \frac{2}{3} 40/2 (0{,}5 \cdot 8{,}057^3 + 12{,}064^3 + 0{,}5 \cdot 14{,}288^3) = \\
        = 46097\,\text{м}^4
    \end{split}
\end{equation*}

Момент инерции потерянной площади относительно поперечной оси в плоскости мидельшпангоута

\begin{equation*}
    \begin{split}
        i'_y = 2 \int_{II}^{I} y x^2 dx = \\
        = 2 l/2 ( 0{,}5 \cdot x_2^2 y_2 + x_3^2 y_3 + 0{,}5 \cdot x_4^2 y_4) = \\
        = 2 \cdot 40/2 (0{,}5 \cdot 60^2 \cdot 8{,}057 + 40^2 \cdot 12{,}064 + 0{,}5 \cdot 20^2 \cdot 14{,}288) = \\
        = 1465984\,\text{м}^4
    \end{split}
\end{equation*}

Момент инерции потерянной площади относительно собственной поперечной оси

\begin{equation*}
    \begin{split}
        i_y = i'_y - s \cdot x_s^2 = \\
        = 1465984 -  928{,}24 \cdot 37{,}34^2 = \\
        = 171711\,\text{м}^4
    \end{split}
\end{equation*}

Далее вычисляются характеристики плавучести и остойчивости судна в поврежденном состоянии

Изменение средней осадки

\begin{equation*}
    \delta T = \frac{v_{\text{в}}}{S - s} = \frac{8293{,}8}{4094{,}2 - 928{,}24} = 2{,}62\,\text{м}
\end{equation*}

Координаты центра тяжести поврежденной ватерлинии

\begin{equation*}
    x_{f1} = \frac{S x_f - s x_s}{S - s} = \frac{4094{,}2 \cdot (-7{,}644) - 928{,}24 \cdot 37{,}34}{4094{,}2 - 928{,}24} = -20{,}83\,\text{м}
\end{equation*}

\begin{equation*}
    y_{f1} = - y_s \frac{s}{S - s} = 0{,}0\,\text{м}
\end{equation*}

Моменты инерции площади поврежденной ватерлинии относительно центральных осей

\begin{equation*}
    I_{x1} = I_x - i_x - s y_s^2 (1 + \frac{s}{S-s}) = 219990 - 46097 - 0 = 173892\,\text{м}^4
\end{equation*}

\begin{equation*}
    \begin{split}
    I_{y1} = I_{yf} - i_y - s (x_s - x_f)^2 (1 + \frac{s}{S-s}) = \\
    =  8179500 - 171711 - 928{,}24 (37{,}34 - (-7{,}644))^2(1 + \frac{928{,}24}{4094{,}2 - 928{,}24}) = \\
    = 5578700\,\text{м}^4
    \end{split}
\end{equation*}

Изменение координат центра величины

\begin{equation*}
    \begin{split}
    \delta x_c = - \frac{v_{\text{в}}}{V} [x_{\text{р}} - x_f + (x_s - x_f)\frac{s}{S-s}] = \\
    = - \frac{8293{,}8}{35360}[37{,}46 - (-7{,}644) + (37{,}34 - (-7{,}644))\frac{928{,}24}{4094{,}2 - 928{,}24}] = \\
    = - 13{,}60\,\text{м}
    \end{split}
\end{equation*}

\begin{equation*}
    \begin{split}
    \delta z_c = - \frac{v_{\text{в}}}{V} [z_{\text{р}} - T - \frac{\delta T}{2}] = \\
    = - \frac{8293{,}8}{35360}[5{,}97 - 10{,}86 - \frac{2{,}62}{2}]\\
    = 1{,}45\,\text{м}
    \end{split}
\end{equation*}

Аварийное возвышение центра величины

\begin{equation*}
    z_{c1} = z_c +\delta z_c = 6{,}04 + 1{,}45 = 7{,}49\,\text{м}
\end{equation*}

Аварийные метацентрические радиусы

\begin{equation*}
    r_1 = \frac{I_{x1}}{V} = \frac{173892}{35360} = 4{,}92\,\text{м}
\end{equation*}

\begin{equation*}
    R_1 = \frac{I_{yf1}}{V} = \frac{5578700}{35360} = 157{,}8\,\text{м}
\end{equation*}

Аварийные метацентрические высоты

\begin{equation*}
    h_1 = z_{c1} + r_1 - z_g = 7{,}49 + 4{,}92 - 12{,}00 = 0{,}41\,\text{м}
\end{equation*}

\begin{equation*}
    H_1 = z_{c1} + R_1 - z_g  = 7{,}49 + 157{,}8 - 12{,}00 = 153{,}3\,\text{м}
\end{equation*}

По приближенной формуле А.~Н.~Крылова вычислим угол дифферента

\begin{equation*}
    \begin{split}
        k_T = \frac{s}{S-s} = \frac{928{,}24}{4094{,}2-928{,}24} = 0{,}293, \\
        k_x = \frac{x_{\text{р}} - x_f}{x_s - x_f} = \frac{37{,}46 - (-7{,}644)}{37{,}34 - (-7{,}644)} = 1{,}003, \\
        \psi = \frac{v_{\text{в}} (x_{\text{р}} - x_f)}{V H_1} (1 + \frac{k_T}{k_x}) = \\
        = \frac{8293{,}8 (37{,}46 - (-7{,}644))}{35630 \cdot 153{,}3} (1 + \frac{0{,}293}{1{,}003}) = 0{,}0892 = 5{,}11\text{\textdegree}.
    \end{split}
\end{equation*}

Аварийные осадка носом и кормой

\begin{equation*}
    \begin{split}
        T_{\text{н}} = T + \delta T + \psi (L/2 - x_{f1}) = \\
        = 10{,}86 + 2{,}62 + 0{,}0892 ( 200/2 - (-20{,}83)) = 24{,}27\,\text{м},
    \end{split}
\end{equation*}

\begin{equation*}
    \begin{split}
        T_{\text{к}} = T + \delta T - \psi (L/2 + x_{f1}) = \\
        = 10{,}86 + 2{,}62 - 0{,}0892 ( 200/2 + (-20{,}83)) = 6{,}42\,\text{м}.
    \end{split}
\end{equation*}

\section{Продольный спуск}

Характеристики судна и спускового устройства:

Спусковой вес судна
\begin{equation*}
    D_{\text{С}} = 0{,}25 \cdot D_{\text{КВЛ}} = 0{,}25 \cdot 35360 = 8840\,\text{т}
\end{equation*}

Абсцисса центра тяжести судна
\begin{equation*}
    x_g = -2{,}50\,\text{м}
\end{equation*}

Аппликата центра тяжести судна
\begin{equation*}
    z_g = 13{,}28\,\text{м}
\end{equation*}

Длина спускового устройства
\begin{equation*}
    L = 0{,}8 \cdot L_{\text{ПП}} = 0{,}8 \cdot 200 = 160\,\text{м}
\end{equation*}

Длина переднего конца полоза
\begin{equation*}
    L_1 = 0{,}5 \cdot L + x_g = 0{,}5 \cdot 160 + (-2{,}50) = 77{,}50\,\text{м}
\end{equation*}

Длина заднего конца полоза
\begin{equation*}
    L_2 = 0{,}5 \cdot L - x_g = 0{,}5 \cdot 160 - (-2{,}50) = 82{,}50\,\text{м}
\end{equation*}

Возвышение киля над зеркалом спусковых дорожек у ЦТ
\begin{equation*}
    C = 0{,}6\,\text{м}
\end{equation*}

Возвышение киля над зеркалом спусковых дорожек у переднего конца полоза
\begin{equation*}
    C_1 = 0{,}7\,\text{м}
\end{equation*}

Возвышение киля над зеркалом спусковых дорожек у заднего конца полоза
\begin{equation*}
    C_2 = 0{,}5\,\text{м}
\end{equation*}

Суммарная ширина полозьев
\begin{equation*}
    b = 2{,}0\,\text{м}
\end{equation*}

Уклон спусковых дорожек
\begin{equation*}
    \beta = 0{,}06
\end{equation*}

Уклон линии киля
\begin{equation*}
    \alpha = \beta + \frac{C_1 - C_2}{L_1 + L_2} = 0{,}06 + \frac{0{,}7 - 0{,}5}{77{,}50 + 82{,}50} = 0{,}05875
\end{equation*}

Расчетная глубина воды на пороге
\begin{equation*}
    T_0 = 0{,}8 \cdot T = 0{,}8 \cdot 3{,}55 = 2{,}84\,\text{м}
\end{equation*}

Длина подводной части дорожек
\begin{equation*}
    \lambda = \frac{T_0}{\beta} = \frac{2{,}84}{0{,}06} = 47{,}3\,\text{м}
\end{equation*}

Удельный вес воды
\begin{equation*}
    \gamma = 1{,}000\,\text{т/м}^3
\end{equation*}

Для спускового веса  судна по кривым элементов теоретического чертежа определяется:

Средняя осадка
\begin{equation*}
    T = 3{,}55\,\text{м}
\end{equation*}

Абсцисса центра величины
\begin{equation*}
    x_c = 0{,}432\,\text{м}
\end{equation*}

Аппликата центра величины
\begin{equation*}
    z_c = 2{,}094\,\text{м}
\end{equation*}

Абсцисса центра тяжести площади ватерлинии
\begin{equation*}
    x_f = 0{,}05\,\text{м}
\end{equation*}

Продольный метацентрический радиус
\begin{equation*}
    R = 610{,}6\,\text{м}
\end{equation*}

Вычисляется продольная метацентрическая высота
\begin{equation*}
    H = R + z_c - z_g = 610{,}6 + 2{,}094 - 13{,}28 = 599{,}4\,\text{м}.
\end{equation*}

Вычисляется дифферент судна
\begin{equation*}
    \psi = \frac{x_g - x_c}{H} = \frac{-2{,}435 - 0{,}432}{599{,}4} = -0{,}0048\,\text{рад}
\end{equation*}

Вычисляется вспомогательная величина --- отстояние центра тяжести площади ватерлинии от центра тяжести судна
\begin{equation*}
    x'_f = x_f - x_g = 0{,}05 - (-2{,}50) = 2{,}55\,\text{м}
\end{equation*}

Углубление передним и задним концами полоза вычисляется по формулам
\begin{equation*}
    \begin{split}
        T_1 = T + C_1 - \psi (L_1 + x'_f) = 3{,}55 + 0{,}7 - 0{,}0048 (77{,}50 - 2{,}55) = 4{,}64\,\text{м} \\
        T_2 = T + C_2 + \psi (L_2 - x'_f) = 3{,}55 + 0{,}5 + 0{,}0048 (82{,}50 - 2{,}55) = 3{,}66\,\text{м}
    \end{split}
\end{equation*}

Осадка $T_2 > T_0$, значит сход судна со стапеля будет сопровождаться соскакиванием.

Приближенное значение баксового давления вычисляется по формуле
\begin{equation*}
    N_{\text{б}} = \frac{D_{\text{С}} H}{L_2 - x'_f} (\alpha + \psi) = \frac{8840 \cdot 599{,}4}{82{,}50 - 2{,}55} (0{,}05875 - 0{,}0048) = 3569\,\text{т}
\end{equation*}

Для водоизмещения
\begin{equation*}
    V_2 = \frac{1}{\rho} (D_{\text{С}} - N_{\text{б}}) = \frac{1}{1{,}000} (8840 - 3569) = 5271\,\text{м}^3
\end{equation*}
определяем значение средней осадки и абсциссу центра тяжести площади ватерлинии
\begin{equation*}
    \begin{split}
        T'_{\text{ср}} = 2{,}35\,\text{м} \\
        x_{f2} = 0{,}699\,\text{м}
    \end{split}
\end{equation*}

Расстояние между центром тяжести площади ватерлинии и центром тяжести судна составляет
\begin{equation*}
    x'_{f2} = x_{f2} - x_g = 0{,}699 - (-2{,}50) = 3{,}20\,\text{м}
\end{equation*}

Тогда ватерлиния всплытия с запасом в 20\% будет проходить на высоте $h$ над килем
\begin{equation*}
    h = 1{,}2 (T'_{\text{ср}} + x'_{f2} \alpha) = 1{,}2 \cdot (2{,}35 + 3{,}20 \cdot 0{,}05875) = 3{,}05\,\text{м}
\end{equation*}

Для определения водоизмещения и центра тяжести погруженной части судна на масштаб Бонжана наносится расчетная ватерлиния по осадкам на перпендикулярах
\begin{equation*}
    \begin{split}
        T_{\text{н}} = h - (L/2 - x_g) \alpha = 3{,}05 - (200 / 2 - (-2{,}50)) \cdot 0{,}05875 = -2{,}98\,\text{м} \\
        T_{\text{к}} = h + (L/2 + x_g) \alpha = 3{,}05 + (200 / 2 + (-2{,}50)) \cdot 0{,}05875 = 8{,}77\,\text{м}
    \end{split}
\end{equation*}

Шаг ватерлиний принимаем равным
\begin{equation*}
    \Delta h = h / n = 3{,}05 / 6 = 0{,}51 \approx 0{,}5 \,\text{м}   
\end{equation*}

Расчет водоизмещения и статического момента для третьей спусковой ватерлинии ($h = 2{,}045\,\text{м}$) приведен в таблице.
Расчет водоизмещения и статического момента для остальных спусковых ватерлиний выполняется аналогично.

% Расчет водоизмещения и статического момента
\begin{table}[H]
    \centering
    \footnotesize
    \caption{Пример расчета водоизмещения }
    \caption*{\textit{
        В столбце II приведены осадки на каждом шпангоуте.
        В столбце III приведены снятые с масштаба Бонжана значения.
        Столбец IV условно назван $V$ --- значение водоизмещения рассчитывается в последней строке столбца.
        Столбец V условно назван $M$ --- значение статического момента  рассчитывается в последней строке столбца.
    }}
    \begin{tikzpicture}
        \matrix (mat) [addenum,
        column 1/.append style={nodes={text width=1.0cm}},
        column 2/.append style={nodes={text width=2.0cm}},
        column 3/.append style={nodes={text width=2.0cm}},
        column 4/.append style={nodes={text width=2.0cm}},
        column 5/.append style={nodes={text width=2.0cm}}
        ] {
        $i_{\text{ШП}}$ & $T$ &  $\Omega$ & $V$ & $M$ \\
        \scriptsize{I}   &    \scriptsize{II}   &   \scriptsize{III}   &   \scriptsize{IV}  &   \scriptsize{V}\\
        0  &  -3,977  &  0,00  &  0,00  &  0,00   \\
        1  &  -2,802  &  0,00  &  0,00  &  0,00   \\
        2  &  -1,627  &  0,00  &  0,00  &  0,00   \\
        3  &  -0,452  &  0,00  &  0,00  &  0,00   \\
        4  &  0,723  &  13,68  &  13,68  &  13,68   \\
        5  &  1,898  &  46,32  &  46,32  &  0,00   \\
        6  &  3,073  &  70,49  &  70,49  &  -70,49   \\
        7  &  4,248  &  74,44  &  74,44  &  -148,88   \\
        8  &  5,423  &  58,48  &  58,48  &  -175,45   \\
        9  &  6,598  &  30,39  &  24,18  &  -96,72   \\
        10  &  7,773  &  0,00  &  0,00  &  0,00   \\
        ---  &  ---  &  ---  &  5751,6  &  -191143   \\
        };
    \end{tikzpicture}
\end{table}

Статический момент относительно плоскости, проходящей через центр тяжести судна, вычисляется по формуле
\begin{equation*}
    M' = M - V x_g = -191143 - 5751{,}6 \cdot (-2{,}50) = -176764\,\text{т}\cdot\text{м}
\end{equation*}

Для построения английской диаграммы спуска необходимы следующие величины в зависимости от пройденного пути $S$:

силы плавучести $\gamma W$;

силы веса $D_{\text{С}}$ (прямая горизонтальная линия);

момента силы плавучести относительно заднего конца полоза $M_W$;

момента силы веса относительно заднего конца полоза $M_D = D_{\text{С}} \cdot L_2$ (прямая горизонтальная линия);

момента силы плавучести относительно порога $M'_{W}$;

момента силы веса относительно порога $M'_{D} = D_{\text{С}} \cdot a$ (наклонная прямая линия).

Расчет величины приведен в таблицах ниже.
Здесь $W = V - v'$, где $v' = 1/2 b \beta l_{\text{п}}^2$ --- потерянный объем плавучести, $l_{\text{п}}$ --- длина спускового полоза, вошедшего в воду и прилегающего к спусковым дорожкам, равная меньшей из величин $\lambda$ или $S$.

% Расчет водоизмещения и статического момента
\begin{table}[H]
    \centering
    \scriptsize
    \caption{Пример расчета английской диаграммы спуска }

    \begin{tikzpicture}
        \matrix (mat) [addenum,
        column 1/.append style={nodes={text width=0.8cm}},
        column 2/.append style={nodes={text width=0.8cm}},
        column 3/.append style={nodes={text width=0.8cm}},
        column 4/.append style={nodes={text width=0.8cm}},
        column 5/.append style={nodes={text width=0.8cm}},
        column 6/.append style={nodes={text width=0.8cm}},
        column 7/.append style={nodes={text width=0.8cm}},
        column 8/.append style={nodes={text width=0.8cm}},
        column 9/.append style={nodes={text width=0.8cm}},
        column 10/.append style={nodes={text width=1.2cm}},
        column 11/.append style={nodes={text width=1.2cm}}
        ] {
        $h$ & $h + c$ &  $\text{II} / \beta$ & $S$ & $V$ & $v'$ & $gW$ & --- & $m'$ & $M'$ & $\gamma m$  \\
        \scriptsize{I}   &    \scriptsize{II}   &   \scriptsize{III}   &   \scriptsize{IV}  &   \scriptsize{V}&   \scriptsize{VI} &   \scriptsize{VII} &   \scriptsize{VIII} &   \scriptsize{IX} &   \scriptsize{X} &   \scriptsize{XI}  \\
        0,545  &  1,14  &  19,1  &  96,6  &  2676,3  &  134,2  &  2542,0  &  -12,45  &  -1671,1  &  -105900  &  -104229  \\
        1,045  &  1,64  &  27,4  &  104,9  &  3560,1  &  134,2  &  3425,9  &  -4,12  &  -552,6  &  -129964  &  -129411  \\
        1,545  &  2,14  &  35,7  &  113,2  &  4580,3  &  134,2  &  4446,1  &  4,22  &  565,9  &  -153931  &  -154497   \\
        2,045  &  2,64  &  44,1  &  121,6  &  5751,6  &  134,2  &  5617,4  &  12,55  &  1684,4  &  -176764  &  -178449 \\
        2,545  &  3,14  &  52,4  &  129,9  &  6972,5  &  134,2  &  6838,3  &  20,88  &  2803,0  &  -199447  &  -202250 \\
        3,045  &  3,64  &  60,7  &  138,2  &  8324,7  &  134,2  &  8190,5  &  29,22  &  3921,5  &  -216933  &  -220855 \\
        };
    \end{tikzpicture}

    \begin{tikzpicture}
        \matrix (mat) [addenum,
        column 1/.append style={nodes={text width=0.8cm}},
        column 2/.append style={nodes={text width=0.8cm}},
        column 3/.append style={nodes={text width=0.8cm}},
        column 4/.append style={nodes={text width=1.0cm}},
        column 5/.append style={nodes={text width=1.2cm}},
        column 6/.append style={nodes={text width=0.8cm}},
        column 7/.append style={nodes={text width=0.8cm}},
        column 8/.append style={nodes={text width=1.2cm}},
        column 9/.append style={nodes={text width=1.2cm}}
        ] {
         $S$ & $x'_c$ &  $\gamma W$ & $L_2 - \text{III}$ & $M_W$ & $a$ & $\text{VII-III}$ & $M'_W$ & $M'd$ \\
        \scriptsize{I}   &    \scriptsize{II}   &   \scriptsize{III}   &   \scriptsize{IV}  &   \scriptsize{V}&   \scriptsize{VI} &   \scriptsize{VII} &   \scriptsize{VIII} &   \scriptsize{IX}  \\
        96,58  &  -41,0  &  2542,0  &  123,5  &  313946  &  -28,2  &  12,79  &  32502  &  -249430   \\
        104,91  &  -37,8  &  3425,9  &  120,3  &  412047  &  -19,9  &  17,89  &  61294  &  -175764   \\
        113,25  &  -34,7  &  4446,1  &  117,2  &  521301  &  -11,5  &  23,20  &  103146  &  -102098   \\
        121,58  &  -31,8  &  5617,4  &  114,3  &  641881  &  -3,2  &  28,55  &  160381  &  -28432   \\
        129,91  &  -29,6  &  6838,3  &  112,1  &  766408  &  5,1  &  34,69  &  237241  &  45234   \\
        138,25  &  -27,0  &  8190,5  &  109,5  &  896573  &  13,5  &  40,42  &  331020  &  118901   \\
        };
    \end{tikzpicture}
\end{table}


Английская диаграмма спуска представлена на рисунке ниже.

\begin{figure}[H]
    \centerline{\incfig{add_11}}
    \caption{Английская диаграмма спуска}
    \caption*{\textit{
        1 --- положение центра тяжести судна по пороге.
        2 --- положение начала всплытия.
        3 --- критическое положение (минимальная разница между $M'_W$ и $M'_D$).
    }}
\end{figure}
